\documentclass[12pt]{article}
\usepackage[utf8]{inputenc}
\usepackage[T1]{fontenc}
\usepackage{hyperref}
\title{Kitan Language Resources for K-12 Curriculum Development}
\author{Compiled by Claude Sonnet 5.5 (claude-sonnet-5-5) for Andr\'as Kornai}
\date{}
\begin{document}
\maketitle

\section{Introduction}
Kitan (Khitan; ISO 639-3 zkt; Glottocode kita1247) is an extinct Mongolic-Khitan language once used in the Liao dynasty. It is known mainly from inscriptions in the Khitan Large and Small scripts. The scripts are only partly deciphered. There are no living speakers and no known heritage community programme. I found no school materials, and no honest K-12 curriculum can use the language ``as much as possible''. At most it could support enrichment units on history, writing systems, and decipherment.

\section{Primary Educational Resources}
I found no primers, textbooks, or school programmes for Khitan.

\section{Linguistic and Reference Materials}
\begin{itemize}
\item Apat\'oczky and Kempf, ``Recent developments on the decipherment of the Khitan Small script'' (\emph{Acta Orientalia Academiae Scientiarum Hungaricae}, 2017): a review with a list of monuments and a bibliography. \url{https://real.mtak.hu/56019/}
\item The same article in the Karoli Gaspar University repository. \url{https://repozitorium.kre.hu/items/0e612016-46ff-4644-b6e4-0854da09ebc3}
\item Wikipedia, List of Khitanologists: scholars and major works, including Janhunen, Wu and Chinggeltei. \url{https://en.wikipedia.org/wiki/List_of_Khitanologists}
\item Kyoto University Hakubi seminar page, ``Current Status of the decipherment of Khitan Script''. \url{https://www.hakubi.kyoto-u.ac.jp/eng/sem/sem_87}
\item Tokyo University of Foreign Studies (AA Institute) news item related to Khitan research. \url{https://www.aa.tufs.ac.jp/en/news/detail/2916/}
\end{itemize}

\section{Religious and Community Resources}
I found none.

\section{Digital and Audio Resources}
I found no audio or software resources. No recordings exist for an extinct language known from inscriptions.

\section{Curriculum Development Context}
There are no speakers, and the phonology and lexicon remain only partly understood. Any school use should be framed as history and linguistics education. Descendant communities, such as the Daur people, are not covered by my search results, so I make no claim about them.
\end{document}